\section{Related work}
\label{sec:related}

\subsection{Origami benchmarks for multimodal models}
\label{sec:related-benchmarks}

\textbf{OrigamiSpace} \citep{origamispace} contributes 350 instances, each carrying a crease
pattern diagram, a compiled flat pattern, the complete folding process and a final folded image,
and defines four tasks over them: pattern prediction, multi-step spatial reasoning, spatial
relationship prediction, and end-to-end crease-pattern code generation. It also provides an
interactive environment and describes reinforcement learning as a possibility explored rather than
a result demonstrated. Its design lesson is that one richly instrumented sample can support
several tasks, which is how 350 instances sustain a benchmark paper.

\textbf{GamiBench} \citep{gamibench} contributes 186 regular and 186 \emph{impossible} crease
patterns paired with folded shapes from six viewpoints, across three visual question-answering
tasks, and introduces two diagnostic metrics: viewpoint consistency and impossible fold selection
rate. Its impossible patterns are the closest the current literature comes to forcing a model to
judge feasibility rather than recognise a shape, and its headline finding is that leading models
struggle even at single-step spatial understanding. It is the benchmark this work is most directly
in conversation with, and the one whose negative examples we can produce without limit, because
generating an illegal fold in our environment is the same operation as generating a legal one.

\textbf{OrigamiBench} \citep{origamibench}, \textbf{FoldingAgent} \citep{foldingagent} and
\textbf{COrigami} \citep{corigami} complete the picture; their metrics appear in the comparison of
Section~\ref{sec:related-gap}.

What all of these share is the gap. In each, the model's proposal is answered by a question rather
than by an execution: it selects among rendered alternatives, or produces an artifact that is
compared to a target. None of them steps a paper-exact engine and asks whether \emph{this} move,
on \emph{this} stack, is something paper could do.

\subsection{Methods that pair a proposer with a simulator}
\label{sec:related-proposer}

\textbf{Learn2Fold} \citep{learn2fold} formulates origami as conditional program induction over a
crease-pattern graph. A language model generates candidate folding programs from text, and a
learned graph-structured world model acts as a differentiable surrogate simulator that predicts
physical feasibility and failure modes before execution, inside a lookahead planning loop. Its
stated key insight is to decouple semantic proposal from physical verification.

We adopt that decomposition and differ on one thing: the verifier. Theirs is learned,
differentiable and approximate, which is what makes it usable for planning and also what makes its
mistakes invisible. Ours is exact, because it is the generator run forwards.
Their scoring is step-level precision, recall and F1 against an expert sequence, together with
edge-level IoU; Section~\ref{sec:scoring-why} shows why that family of measures is unsafe in this
task, since a model that finds a \emph{shorter} correct sequence is marked wrong by it. That is a
finding about scoring in this setting rather than a criticism of their results.

The general precedent for a (((((proposer paired with a verifier))))))) - ----what's this??? what's
pur proposer
is older than any of this work: a
network that proposes and a search that checks is the structure behind AlphaGo \citep{alphago}, and
the reason the
combination is stronger than either half.

\subsection{Computational origami: the theory under the action space}
\label{sec:related-theory}

The benchmark's central claim, that geometry does not determine a folded state, is Bern and
Hayes's: deciding the overlap order of a flat folding is NP hard even given a valid
mountain-valley assignment \citep{bern-hayes}. The action space is Arkin et al.'s: simple folds,
in the one-layer, some-layers and all-layers variants, with finite or infinite fold lines
\citep{arkin-map}. Their complexity has been mapped since, including simple foldability's hardness
\citep{akitaya-hard} and the infinite all-layers case \citep{akitaya-infinite}. The solver whose
constraint vocabulary we borrow for terminal states is Flat-Folder \citep{akitaya-flatfolder}.

Two results shape what the task actually is. \citet{continuous} showed that any well-behaved
folded state is reachable by \emph{some} continuous motion; the question is therefore never
whether a path exists, but whether one exists through a finite number of the discrete operations a
hand or a machine can perform. And counting foldings has no closed form even for a
one-dimensional strip, a problem open since Lucas attributed it to Lemoine in 1891
\citep{lucas,oeis-a000136}; this is the answer to ``why not enumerate the states'', and it is a
fact about the problem rather than about any implementation.

Finally, the problem itself is not new: generating folding sequences from crease patterns was
named and attacked symbolically by \citet{akitaya-cp2seq}, by the same group whose tools this
field now depends on. What is new is neither the problem nor the theory, but the observation that
this problem supplies an exact verifier and can therefore serve as an evaluation substrate for
models that reason with tools.

\subsection{The structural gap}
\label{sec:related-gap}

The two closest benchmarks are worth setting beside this one directly, because a reader who knows
them will ask what is left to do.

\begin{table}[h]
\caption{The two closest benchmarks beside this one. The informative rows are the last three.}
\label{tab:closest}
\begin{center}
\small
\begin{tabular}{p{0.15\textwidth}p{0.21\textwidth}p{0.23\textwidth}p{0.27\textwidth}}
 &\multicolumn{1}{c}{\bf GAMIBENCH} &\multicolumn{1}{c}{\bf ORIGAMISPACE} &\multicolumn{1}{c}{\bf CP2SEQ (OURS)} \\
\hline \\
Model produces & A choice among rendered options & A predicted pattern, a relation, or crease-pattern code & A sequence of folds, one step at a time \\[4pt]
Judged by & Agreement with the stored answer & Similarity to a stored target, or compilation & Execution of each step, then replay of the whole \\[4pt]
Judge is & A stored key & A comparison against a reference artifact & The fold engine itself, run forwards \\[4pt]
Judge can be wrong & Not applicable; no execution & Where similarity stands in for correctness & Only if the model of paper is wrong, which replay cannot detect (Sec.~\ref{sec:cp2seq-verify}) \\[4pt]
Wrong answers & 186 authored impossible patterns & Not a designed axis & Unlimited: an illegal fold is produced by asking for one \\[4pt]
Ground truth & Annotated per instance & Annotated per instance & Constructed; recorded as the fold happens \\[4pt]
Non-unique answers & One key per question & Compared against one reference & Replayed and compared up to a declared symmetry group (Sec.~\ref{sec:scoring-what}) \\[4pt]
Scores search effort & No & No & Yes: queries to solution, over all attempts (Sec.~\ref{sec:scoring-metrics}) \\
\end{tabular}
\end{center}
\end{table}

Read across the bottom three rows rather than the top. GamiBench's impossible patterns are the
closest the literature comes to forcing a judgement of feasibility rather than a recognition of
shape, and they are authored, which is why there are 186 of them; here the same object is a
by-product of generation and there is no limit on it. OrigamiSpace's instrumentation per instance
is richer than ours and is the reason 350 instances sustain a benchmark paper; what it does not
have is a judge that can rule on a step the authors never anticipated. Neither ranks among the
several valid folded states of one crease pattern, because both compare against a stored answer,
and a stored answer cannot represent a set.

None of this is a defect in either benchmark. Both were built to measure whether a model
understands folded geometry, and they measure it. The claim here is narrower: that measuring
whether a model can \emph{search} for a folding sequence requires a judge that executes, and that
origami supplies one.

Collecting the metrics used across this literature makes one absence visible. Every metric in use
compares a produced artifact against a target: geometric or semantic similarity to a reference,
precision and recall against an expert sequence, IoU over affected edges, compilation success,
human preference. \textbf{None ranks among the several valid folded states of a single crease
pattern, and none scores an agent by how much verification it required.} The absence is structural
rather than accidental, because computing any of those metrics requires a target to compare
against, and the interesting question here has more than one right answer.

